\section{模型评价与推广}
\subsection{模型优势}
我们在四问中采用相同的供需平衡和储电量递推关系，并根据可用预报、合同调整条件与电价描述不同购电任务；储能在有余电时充电、存在缺口时放电，从而明确计划量、补购量和费用之间的关系。

在含最坏损失限制的实验对照下，我们的消融结果支持净负荷保护、光伏组合预测和日内合同调整的作用。去掉日前保护后，两种电价下的费用分别增加2.01\%和2.21\%；去掉光伏组合后分别增加1.78\%和1.97\%。保持该对照的预测与保护安排时，允许合同调整分别减少23.39万元和26.95万元费用。

\subsection{适用条件与局限}
我们在正式方案中保留净增益均值门槛并取消最坏情景损失限制。相对含该限制的对照，固定与波动电价下的费用分别减少7.15万元和5.98万元，紧急购电量也降低，日费用与紧急电量的最大值及95\%分位数均下降。

2025年数据已参与此前配置选择，本次消融与敏感性分析属于同一年度的回顾性比较。模型的适用前提是近期误差仍有代表性、紧急购电能够保障供电、储能效率近似恒定；跨年度效果需用新增数据检验。

\subsection{模型改进与推广}
在含最坏损失限制的对照及所设扰动范围内，我们发现设备效率变化引起的费用变化最大：单向充放电效率由0.90降至0.85时，固定与波动电价下的费用分别增加3.54\%和3.21\%。推广时，应先根据设备实测数据估计效率随储电状态、充放电功率的变化，再加入循环寿命成本、功率爬坡和联络线容量限制。

针对风险限制的保守代价，我们可用新增年度数据比较不同历史误差窗口与损失预算下的实际费用、紧急购电和不利日期，选择与供电要求相适应的调整规则；气温、节假日等信息也可用于预测，并同时检验预测误差和购电费用。
